\section{Hybrid CNN + Self-Attention: Theoretical Background}
This section establishes the theory behind the hybrid architecture that combines a convolutional front-end with a Transformer encoder. The convolutional half reuses every concept from the 1D-CNN section, so the focus here is on what the convolution alone cannot do, and how self-attention complements it.

\subsection{The Locality Limitation of Pure CNNs}
A convolutional layer sees only a fixed window of size $K$ at every position. To enlarge its \textbf{receptive field} --- the span of input time steps that influence a single output --- the network must either stack many layers (slow growth, linear in depth) or apply pooling (loss of resolution). For a stack of $L$ layers with kernel size $K$ and a single pooling factor of 2 between each pair, the receptive field grows as roughly $K \cdot 2^{L-1}$. To cover a 500-point anomaly with $K = 7$, the network needs at least 7 stacked layers, which is far deeper than a typical 1D-CNN baseline.

Equally important, pure convolutions enforce \textbf{translation equivariance} but cannot reason about \textbf{long-range dependencies} in a flexible way. The convolution at time $t$ has no mechanism to selectively attend to a single distant pattern at $t' \gg t$ while ignoring everything in between --- it sees a uniformly weighted local patch. For classes like \texttt{M} (two peaks separated by a trough), this matters: distinguishing an \texttt{M} from two separate \texttt{bell} events requires reasoning about the relationship between events 300 points apart, not just summing local features.

\subsection{Scaled Dot-Product Attention}
Self-attention solves both problems by letting every position attend to every other position with a \textit{data-dependent} weight. Given a sequence of $T$ feature vectors $\mathbf{X} \in \mathbb{R}^{T \times d_{\text{model}}}$, three learned linear projections produce a \textit{query}, a \textit{key}, and a \textit{value} representation:
\begin{equation}
    \mathbf{Q} = \mathbf{X} \mathbf{W}_{Q}, \quad
    \mathbf{K} = \mathbf{X} \mathbf{W}_{K}, \quad
    \mathbf{V} = \mathbf{X} \mathbf{W}_{V}
\end{equation}
where $\mathbf{W}_{Q}, \mathbf{W}_{K}, \mathbf{W}_{V} \in \mathbb{R}^{d_{\text{model}} \times d_{k}}$. The attention output is then the value matrix re-weighted by a softmax-normalised query-key similarity:
\begin{equation}
    \text{Attention}(\mathbf{Q}, \mathbf{K}, \mathbf{V}) = \text{softmax}\!\left( \frac{\mathbf{Q} \mathbf{K}^{\top}}{\sqrt{d_{k}}} \right) \mathbf{V}
\end{equation}
The intuition is the soft-lookup analogy. Each row of $\mathbf{Q} \mathbf{K}^{\top}$ measures how relevant every key (position) is to a given query (position); the softmax converts these scores into a probability distribution; the matrix product against $\mathbf{V}$ produces a weighted average of value vectors. The output at position $t$ is therefore a context-dependent mixture of features from every other position in the sequence, with the mixture weights learned end-to-end.

The $\sqrt{d_{k}}$ scaling is not cosmetic. For large $d_{k}$, the dot products $\mathbf{q}_{t} \cdot \mathbf{k}_{s}$ have variance proportional to $d_{k}$, pushing the softmax into a saturated regime where the gradient vanishes. Dividing by $\sqrt{d_{k}}$ keeps the pre-softmax logits unit-variance and preserves the gradient signal during backpropagation.

\subsection{Multi-Head Attention}
A single attention head can only learn one type of relationship. Vaswani et al.\ (2017) address this by running $h$ attention operations in parallel, each with its own projections, and concatenating their outputs:
\begin{equation}
    \text{MultiHead}(\mathbf{X}) = [\text{head}_{1}; \ldots; \text{head}_{h}] \mathbf{W}_{O}, \quad
    \text{head}_{i} = \text{Attention}(\mathbf{X} \mathbf{W}_{Q}^{(i)}, \mathbf{X} \mathbf{W}_{K}^{(i)}, \mathbf{X} \mathbf{W}_{V}^{(i)})
\end{equation}
Each head dimensionally projects to $d_{k} = d_{\text{model}} / h$, so the total parameter and compute cost matches a single head at full dimension. Different heads typically specialise: in our setting, one head may track the leading edge of an anomaly while another tracks the trailing edge, and a third may attend to the noise baseline for context.

\subsection{Sinusoidal Positional Encoding}
Attention is, by construction, \textbf{permutation-equivariant}: shuffling the input positions produces a correspondingly shuffled output. For a sequence model, this is catastrophic --- time order matters. The standard fix is to add a positional code to the input embeddings:
\begin{equation}
    \text{PE}_{t, 2i} = \sin\!\left( \frac{t}{10000^{2i / d_{\text{model}}}} \right), \quad
    \text{PE}_{t, 2i+1} = \cos\!\left( \frac{t}{10000^{2i / d_{\text{model}}}} \right)
\end{equation}
The sinusoidal choice has two useful properties. First, every relative offset $\Delta t$ can be expressed as a linear transformation of the positional code, allowing the model to easily learn shift-invariant comparisons. Second, the encoding generalises to sequences longer than those seen during training, since the formulas are defined for any positive $t$.

\subsection{The Transformer Encoder Layer}
A single encoder layer wraps multi-head attention with two residual sub-layers and a position-wise feed-forward network:
\begin{align}
    \mathbf{X}' &= \text{LN}(\mathbf{X} + \text{MultiHead}(\text{LN}(\mathbf{X}))) \\
    \mathbf{X}'' &= \text{LN}(\mathbf{X}' + \text{FFN}(\text{LN}(\mathbf{X}')))
\end{align}
where $\text{FFN}(\mathbf{x}) = \mathbf{W}_{2} \, \text{GELU}(\mathbf{W}_{1} \mathbf{x} + \mathbf{b}_{1}) + \mathbf{b}_{2}$ is applied independently to every position. We use the \textbf{pre-norm} variant (LayerNorm before each sub-layer, as shown above) rather than the original \textbf{post-norm} of Vaswani et al.\ for one practical reason: pre-norm Transformers are dramatically more stable during the first few hundred training steps and can be trained without a learning-rate warm-up schedule.

\subsection{The CLS Token Aggregation Strategy}
After $L$ encoder layers, the network holds a sequence of $T$ contextualised feature vectors. A classifier expects a single fixed-size vector. There are two standard ways to aggregate: global pooling across positions, or the \textbf{CLS token} trick borrowed from BERT (Devlin et al., 2018). We use the latter: a single learnable vector $\mathbf{c} \in \mathbb{R}^{d_{\text{model}}}$ is prepended to the sequence before the first encoder layer. After the final layer, only the position originally occupied by $\mathbf{c}$ is fed to the classifier head.

The CLS approach is elegant because it lets the network \textit{learn how to summarise}. The attention mechanism in every encoder layer is free to aggregate any subset of positions into the CLS slot, with the mixture weights driven by the classification loss. By contrast, a global average pool gives every time step equal weight regardless of whether it contains an anomaly or pure background noise.

\newpage
